\section{Open Questions}

This replication verified the exact analytic backbone but leaves five substantive
questions open for follow-up work. Each is grounded in what was (and was not) tested
here, with concrete next steps that would close the gap.

\subsection*{Q1. Extension to non-diagonal / non-standard magic states}
\textbf{Basis.} All exact $\xi$ values obtained here ($\xi(T) = 1.171573$,
$\xi(\mathrm{CCZ}) = 16/9$, closed-form $\xi(R(\theta))$) are for \emph{diagonal}
magic states in the standard qubit Clifford + $T$ setting. Prop.~2 ($\xi = 1/F$)
additionally requires the state be a Clifford magic state. Whether the same rank
bounds hold for $Y$-basis rotations, controlled-phase-of-arbitrary-angle gates,
or qutrit magic states is not derived by the paper and not tested by us.

\textbf{Next steps.} (a) Extend the brute-force maximizer in
\texttt{src/verify\_extent.py} to a family of magic states parameterized by rotation
axis and angle, checking whether $\xi = 1/F$ still holds when the Clifford-magic-state
condition is relaxed. (b) For qudit primes $p = 3, 5$, enumerate the generalized-Pauli
stabilizer set (feasible for small $n$ via BFS) and repeat the CCZ-analog search.
(c) Compare with the qudit magic-state literature (Howard \& Vala, 2012).

\subsection*{Q2. Tightness of Bravyi--Smith exact-rank bounds $\chi(T^m)$}
\textbf{Basis.} The paper cites $\chi(T^m) = 1, 2, 3, 4, 5, 7, 12, \ldots$ from Ref.~[14]
with the note that the search required heavy compute. Whether $\chi(T^7) = 12$ is
truly minimal, or a lower rank exists via a non-obvious stabilizer combination, is
not resolved by any exhaustive proof visible in the paper. Our C8 was not attempted.

\textbf{Next steps.} (a) Frame rank-minimization as compressed-sensing / low-rank
matrix recovery over the stabilizer dictionary and use IRLS or basis-pursuit as a
lower-bound heuristic. (b) SAT/SMT-encode ``exists a rank-$\leq k$ decomposition of
$|T^m\rangle$?'' for $m = 7, 8$; UNSAT certificates yield tight lower bounds.
(c) Request the raw enumeration data from the Ref.~[14] authors for an independent
audit.

\subsection*{Q3. GPU acceleration of the sum-over-Cliffords sampler}
\textbf{Basis.} Our C6 verification ran at $t \leq 10$ (max err $1.7 \times 10^{-15}$
exact, $0.01$--$0.08$ sampled at $k \sim 10^2$), four orders of magnitude below the
paper's demonstrated $t = 40$--$64$ regime. The estimator's bias and $O(\delta^{-2}
\prod \xi_j)$ variance scaling under a GPU tableau backend are not tested.

\textbf{Next steps.} (a) Port \texttt{src/verify\_soc\_sim.py} to a JAX or CUDA-Q
bit-packed symplectic-tableau backend; benchmark samples/sec vs.\ CPU \texttt{numpy}.
(b) Run a variance-vs-$k$ sweep at $t = 20, 30, 40$ on GPU and confirm empirical
variance matches $\delta^{-2} \prod_j \xi(V_j)$. (c) Cross-check against Google's
\texttt{stim} (Gidney, 2021).

\subsection*{Q4. Composition with magic-state distillation circuits}
\textbf{Basis.} The paper's $\xi$ machinery treats each non-Clifford gate as a fixed
magic-state input with known $\xi$. In a full fault-tolerant stack, $T$-states are
themselves the \emph{output} of distillation subcircuits (15-to-1 Reed--Muller,
Reichardt, Fowler--Devitt); the effective $\xi$ of a distilled state depends on the
protocol's residual noise. This coupling is not discussed.

\textbf{Next steps.} (a) For the 15-to-1 Reed--Muller circuit, compute $\xi$ of the
distilled $T$-state as a function of raw-$T$ infidelity $\varepsilon$, and confirm
convergence to $\xi(T) = 1.1716$ as $\varepsilon \to 0$. (b) Integrate with a
surface-code decoder simulator (PyMatching + \texttt{stim}) for end-to-end
distilled-$\xi$ vs.\ code distance. (c) Publish the $\xi(\varepsilon)$ curve as a
resource-estimation input for FTQC compilers.

\subsection*{Q5. Practical classical baseline for near-term variational algorithms}
\textbf{Basis.} The paper demonstrates $50$-qubit QAOA and Hidden-Shift as headline
runs but does not sweep the $(n, t, \text{depth}, \text{ansatz})$ envelope inside
which the framework beats state-vector, tensor-network (MPS/PEPS), or neural-network
quantum-state simulators. VQE / UCCSD chemistry ansatze with structured $T$-count
reductions are not addressed.

\textbf{Next steps.} (a) Small chemistry benchmark (H$_2$, LiH, BeH$_2$ with UCCSD
on $6$--$14$ qubits): compare low-rank stabilizer wall-clock vs.\ exact statevector
vs.\ MPS at $\delta = 1$~mHa. (b) QAOA-MaxCut $p = 1, 2$ with $n = 30$--$60$: run
the paper's simulator alongside \texttt{cotengra}/\texttt{quimb} and publish a Pareto
frontier $(n, t, \text{wall-clock}, \text{memory})$. (c) Document the
$(n_{\max}, t_{\max})$ envelope inside which low-rank stabilizer beats both
statevector and tensor networks --- a prerequisite for honestly benchmarking
hardware ``quantum advantage'' claims.
